\begin{abstract}
Estimation errors in the reproduction number are frequently attributed to sample size. Here we show that the return on increasing case counts depends fundamentally on the inferential target, distinguishing the conditional instantaneous reproduction number $R_t$ from the cross-period mean $\bar{R}$. Under negative-binomial transmission and binomial reporting, the estimation error of $R_t$ decays as $1/m$ with expected reporting scale $m$, attaining the Cram\'er--Rao bound. For single-period estimators of $\bar{R}$ that are conditionally unbiased for $R_t$, however, the estimation error retains an irreducible variance floor $v_R$ arising from environmental stochasticity, with reporting heterogeneity introducing an additional component $v_H$ into our estimator's variance when reporting probabilities fluctuate across periods. This irreducibility is specific to expanding single-period case volume and does not preclude variance reduction through accumulating independent observation periods. We derive the variance balance scale $m_\times = A/v_R$, which characterises the point of diminishing sample returns without prescribing a minimum operational threshold, and reduce surveillance optimisation to three levers: sample volume, period count, and reporting standardisation. Analysing 10,303 valid state-weeks of influenza hospitalisations from the US National Healthcare Safety Network (2020--2026), we find that early-epidemic squared error scales affinely with inverse sample size (intercept $a = 0.0991$, 95\% CI 0.0837--0.1165; slope $b = 2.243$, 1.494--3.007; ratio $b/a = 22.6$, 13.5--34.9). Theoretical analysis indicates that expanding single-period sample volume yields diminishing returns; empirical observations are compatible with this error structure, though current data cannot separate a strict non-zero asymptotic plateau from slow decay or uniquely attribute the fitted intercept to environmental variance. Reporting adjustments account for $<1\%$ of the intercept excess. Surveillance design should distinguish outbreak detection from estimation of mean transmission, and balance additional case capture against temporal replication, spatial pooling, and reporting standardisation according to the inferential target and dependence structure.
\end{abstract}
